% ---
% title: Handwritten Medicine MCQs
% date: 2026-09
% category:
% nodes: tools 
% links: 
% ---
\documentclass{article}
\usepackage[utf8]{inputenc}
\usepackage{amsmath, amssymb, amsthm}
\usepackage{geometry}
\usepackage{graphicx}
\usepackage{hyperref}

\geometry{a4paper, margin=1in}

\title{Handwritten Medicine MCQs}
\date{August 2026}
\author{Daksh Mehta}

\begin{document}
	
Here are some multiple choice and short answer questions written for UKMLA style examinations. \href{mailto:daksh1111@proton.me}{Email me} if there are any discrepancies, and \href{handwritten-medicine-mcq-answers.html}{click here for the answers}. I will continue to add more on a very irregular basis with varying levels of difficulty. These questions are based on concepts I personally haven't seen tested before in other question banks I've used. I have no idea whether they would come up on a real exam but I find them interesting nonetheless. 

\tableofcontents

\section{Hypersensitivity reactions in SLE}

A 28-year-old woman with systemic lupus erythematosus presents to the emergency department with a painful, swollen left leg. Laboratory studies show a prolonged activated partial thromboplastin time with a normal bleeding time. Testing reveals antibodies against beta 2-glycoprotein I and cardiolipin. She has previously experienced two unexplained miscarriages. What type of hypersensitivity reaction is the underlying cause for the patient's current presentation?

a. Type I  
b. Type II  
c. Type III  
d. Type IV

\section{Retinal haemorrhages in the neonate}

A 10-day-old boy is brought to the emergency department with irritability and poor feeding. Examination shows bilateral retinal haemorrhages and several ecchymoses on the trunk. There is also persistent bleeding from the umbilical stump. He was born at term following an uncomplicated pregnancy and is exclusively breastfed. Which of the following is the most likely explanation for his bleeding?

a. Deficiency of factor VIII \\
b. Maternal antibodies against platelet glycoproteins \\
c. Vitamin K deficiency \\ 
d. Disseminated intravascular coagulation \\
e. Shaken baby syndrome


	
\end{document}